\documentclass[12pt]{article}
\usepackage[utf8]{inputenc}
\usepackage[T1]{fontenc}
\usepackage{hyperref}
\title{Northern Tujia Language Resources for K-12 Curriculum Development}
\author{Compiled by Claude Sonnet 5 (claude-sonnet-5) for Andr\'as Kornai}
\date{}
\begin{document}
\maketitle
\section{Introduction}
Northern Tujia (ISO 639-3 tji) is a Sino-Tibetan language of Xiangxi, China, with an estimated 70,000 speakers as of 2005. It is seriously endangered and has no established written form. The main scholarly resource found is the Brassett grammar, in English.
\section{Primary Educational Resources}
No primers, textbooks or school programmes were found in the search results.
\section{Linguistic and Reference Materials}
\begin{itemize}
\item Review of the Brassetts' English grammar of Northern Tujia, based on 18 months of fieldwork in Jishou. It has about 200 pages, a glossary of over 1,500 items and 3 annotated texts: \url{https://linguistlist.org/issues/18/311/}
\item LINGUIST List issue 17.288, also relating to Tujia: \url{https://linguistlist.org/issues/17/288}
\item Wikipedia article on the Tujia language: \url{https://en.wikipedia.org/wiki/Tujia_language}
\end{itemize}
\section{Religious and Community Resources}
\begin{itemize}
\item Tujia language and culture page (not checked): \url{https://tujiayu.hu60.cn/elng/elintro.html}
\end{itemize}
\section{Digital and Audio Resources}
\begin{itemize}
\item Endangered Languages Project entry for Northern Tujia: \url{https://www.endangeredlanguages.com/lang/4225?hl=es}
\item ELDP documentation of Southern Tujia, a different, mutually unintelligible variety, listed for context only: \url{https://eldp.access.preservica.com/dk0260/}
\end{itemize}
\section{Curriculum Development Context}
No school use is known. A curriculum would need an agreed writing system and speaker involvement. No standalone dictionary was found, only the glossary inside the grammar.
\end{document}